% Source PDF pages 76--78; manual pages 2-47--2-49.
\subsection{Flight Path Angle Rates}\label{sec:flight-path-rates}

The function \texttt{compute\_rates} evaluates the following flight-path-angle rate
output variables:
\begin{longtable}{@{}>{\ttfamily}p{0.17\textwidth}p{0.19\textwidth}p{0.56\textwidth}@{}}
  psigcdt & $\dot\psi_{gc}$ & Geocentric heading rate.\\
  psigddt & $\dot\psi_{gd}$ & Geodetic heading rate.\\
  gamgcdt & $\dot\gamma_{gc}$ & Geocentric vertical flight-path-angle rate.\\
  gamgddt & $\dot\gamma_{gd}$ & Geodetic vertical flight-path-angle rate.\\
\end{longtable}

The rates are derived from the angular velocity between the velocity coordinate
system defined by \cref{eq:velocity-x-from-geodetic,eq:velocity-y-from-geodetic,eq:velocity-z-from-geodetic}
and the geocentric or geodetic coordinate system. The derivation below is for the
geodetic rates; the geocentric derivation is identical after replacing $gd$ with
$gc$.

The angular velocity of the velocity coordinate system is
\begin{equation}
  {}_{\oplus}\vec\omega_V
  ={}_{\oplus}\vec\omega_{gd}
   +{}_{gd}\vec\omega_V,
  \label{eq:velocity-frame-angular-velocity-sum}
\end{equation}
the sum of the angular velocity between ECFC and geodetic coordinates and the
angular velocity between geodetic and velocity coordinates. From the coordinate
system definitions,
\begin{equation}
  {}_{\oplus}\vec\omega_{gd}
  =\dot\lambda\,\hat z_{\oplus}
   -\dot\delta_{gd}\,\hat y_{gd},
  \label{eq:ecfc-to-geodetic-angular-velocity}
\end{equation}
and
\begin{equation}
  {}_{gd}\vec\omega_V
  =\dot\psi_{gd}\,\hat z_{gd}
   +\dot\gamma_{gd}\,\hat y_V.
  \label{eq:geodetic-to-velocity-angular-velocity}
\end{equation}
Therefore,
\begin{equation}
  {}_{\oplus}\vec\omega_V
  =\dot\lambda\,\hat z_{\oplus}
   -\dot\delta_{gd}\,\hat y_{gd}
   +\dot\psi_{gd}\,\hat z_{gd}
   +\dot\gamma_{gd}\,\hat y_V.
  \label{eq:velocity-frame-angular-velocity-expanded}
\end{equation}
The component in the $\hat y_V$ direction is
\begin{equation}
  {}_{\oplus}\vec\omega_V\mathbin{\cdot}\hat y_V
  =\left(
      \dot\lambda\,\hat z_{\oplus}
      -\dot\delta_{gd}\,\hat y_{gd}
    \right)\mathbin{\cdot}\hat y_V
   +\dot\gamma_{gd},
  \label{eq:velocity-angular-rate-y-component}
\end{equation}
because
\begin{equation}
  \dot\psi_{gd}\,\hat z_{gd}\mathbin{\cdot}\hat y_V=0.
  \label{eq:heading-term-orthogonal-yv}
\end{equation}
Using the heading-angle relationship in \cref{eq:velocity-y-from-geodetic},
\cref{eq:velocity-angular-rate-y-component} becomes
\begin{equation}
  {}_{\oplus}\vec\omega_V\mathbin{\cdot}\hat y_V
  =\dot\lambda\,\hat z_{\oplus}\mathbin{\cdot}\hat y_V
   -\dot\delta_{gd}\cos\psi_{gd}
   +\dot\gamma_{gd}.
  \label{eq:velocity-angular-rate-y-simplified}
\end{equation}
The component of \cref{eq:velocity-frame-angular-velocity-expanded} in the
$\hat z_V$ direction is
\begin{equation}
  {}_{\oplus}\vec\omega_V\mathbin{\cdot}\hat z_V
  =\left(
      \dot\lambda\,\hat z_{\oplus}
      -\dot\delta_{gd}\,\hat y_{gd}
      +\dot\psi_{gd}\,\hat z_{gd}
    \right)\mathbin{\cdot}\hat z_V.
  \label{eq:velocity-angular-rate-z-component}
\end{equation}
Using \cref{eq:velocity-z-from-geodetic} simplifies this expression to
\begin{equation}
  {}_{\oplus}\vec\omega_V\mathbin{\cdot}\hat z_V
  =\dot\lambda\,\hat z_{\oplus}\mathbin{\cdot}\hat z_V
   -\dot\delta_{gd}\sin\gamma_{gd}\sin\psi_{gd}
   +\dot\psi_{gd}\cos\gamma_{gd}.
  \label{eq:velocity-angular-rate-z-simplified}
\end{equation}

An independent expression for the angular velocity between the ECFC and velocity
systems follows by differentiating the velocity in the velocity frame:
\begin{equation}
  \vec a_{\oplus}
  =\frac{d\lVert\vec V_{\oplus}\rVert}{dt}\,\hat x_V
   +{}_{\oplus}\vec\omega_V\times\vec V_{\oplus},
  \label{eq:velocity-derivative-in-rotating-velocity-frame}
\end{equation}
where
\begin{equation}
  \frac{d\lVert\vec V_{\oplus}\rVert}{dt}
  =\frac{\vec a_{\oplus}\mathbin{\cdot}\vec V_{\oplus}}
         {\lVert\vec V_{\oplus}\rVert}
  =\vec a_{\oplus}\mathbin{\cdot}\hat x_V.
  \label{eq:speed-rate-from-acceleration}
\end{equation}
Because $\vec V_{\oplus}$ is aligned with $\hat x_V$, the cross product becomes
\begin{equation}
  \begin{aligned}
  {}_{\oplus}\vec\omega_V\times\vec V_{\oplus}
  &=\begin{vmatrix}
    \hat x_V&\hat y_V&\hat z_V\\
    \omega_{Vx}&\omega_{Vy}&\omega_{Vz}\\
    \lVert\vec V_{\oplus}\rVert&0&0
  \end{vmatrix}\\
  &=\lVert\vec V_{\oplus}\rVert
    \left[
      \left({}_{\oplus}\vec\omega_V\mathbin{\cdot}\hat z_V\right)\hat y_V
      -\left({}_{\oplus}\vec\omega_V\mathbin{\cdot}\hat y_V\right)\hat z_V
    \right].
  \end{aligned}
  \label{eq:velocity-frame-cross-product}
\end{equation}
The quantities $\omega_{Vx}$, $\omega_{Vy}$, and $\omega_{Vz}$ are the components
of ${}_{\oplus}\vec\omega_V$ in velocity coordinates. Combining
\cref{eq:velocity-derivative-in-rotating-velocity-frame,eq:speed-rate-from-acceleration,eq:velocity-frame-cross-product}
gives
\begin{equation}
  \begin{aligned}
  \vec a_{\oplus}={}&
  \left(\vec a_{\oplus}\mathbin{\cdot}\hat x_V\right)\hat x_V\\
  &+\lVert\vec V_{\oplus}\rVert
    \left[
      \left({}_{\oplus}\vec\omega_V\mathbin{\cdot}\hat z_V\right)\hat y_V
      -\left({}_{\oplus}\vec\omega_V\mathbin{\cdot}\hat y_V\right)\hat z_V
    \right].
  \end{aligned}
  \label{eq:acceleration-in-velocity-frame}
\end{equation}
Its $\hat y_V$ component is
\begin{equation}
  \vec a_{\oplus}\mathbin{\cdot}\hat y_V
  =\lVert\vec V_{\oplus}\rVert
    \left({}_{\oplus}\vec\omega_V\mathbin{\cdot}\hat z_V\right),
  \label{eq:acceleration-yv-component}
\end{equation}
so
\begin{equation}
  {}_{\oplus}\vec\omega_V\mathbin{\cdot}\hat z_V
  =\frac{\vec a_{\oplus}\mathbin{\cdot}\hat y_V}
         {\lVert\vec V_{\oplus}\rVert}.
  \label{eq:angular-rate-z-from-acceleration}
\end{equation}
Similarly,
\begin{equation}
  \vec a_{\oplus}\mathbin{\cdot}\hat z_V
  =-\lVert\vec V_{\oplus}\rVert
    \left({}_{\oplus}\vec\omega_V\mathbin{\cdot}\hat y_V\right),
  \label{eq:acceleration-zv-component}
\end{equation}
and therefore
\begin{equation}
  {}_{\oplus}\vec\omega_V\mathbin{\cdot}\hat y_V
  =-\frac{\vec a_{\oplus}\mathbin{\cdot}\hat z_V}
          {\lVert\vec V_{\oplus}\rVert}.
  \label{eq:angular-rate-y-from-acceleration}
\end{equation}
Equating \cref{eq:velocity-angular-rate-y-simplified} and
\cref{eq:angular-rate-y-from-acceleration} gives the geodetic vertical
flight-path-angle rate:
\begin{equation}
  \dot\gamma_{gd}
  =-\frac{\vec a_{\oplus}\mathbin{\cdot}\hat z_V}
          {\lVert\vec V_{\oplus}\rVert}
   -\dot\lambda\,\hat z_{\oplus}\mathbin{\cdot}\hat y_V
   +\dot\delta_{gd}\cos\psi_{gd}.
  \label{eq:geodetic-flight-path-angle-rate}
\end{equation}
Equating \cref{eq:velocity-angular-rate-z-simplified} and
\cref{eq:angular-rate-z-from-acceleration} gives the geodetic heading rate:
\begin{equation}
  \dot\psi_{gd}
  =\frac{
    \dfrac{\vec a_{\oplus}\mathbin{\cdot}\hat y_V}
          {\lVert\vec V_{\oplus}\rVert}
    -\dot\lambda\,\hat z_{\oplus}\mathbin{\cdot}\hat z_V
    +\dot\delta_{gd}\sin\gamma_{gd}\sin\psi_{gd}
  }{\cos\gamma_{gd}}.
  \label{eq:geodetic-heading-rate}
\end{equation}
The velocity-frame unit vectors $\hat y_V$ and $\hat z_V$ are given by
\cref{eq:velocity-y-from-geodetic,eq:velocity-z-from-geodetic}.
